% Reconstructed from the 14 August 2026 FINAL PDF.
\documentclass[11pt]{article}
\usepackage[margin=1.1in]{geometry}
\usepackage{amsmath,amssymb,amsthm}
\usepackage{enumitem}
\usepackage[colorlinks=true,linkcolor=blue,citecolor=blue,urlcolor=blue]{hyperref}

\newcommand{\pP}{\textbf{[P]}}
\newcommand{\pC}{\textbf{[C]}}
\newcommand{\pO}{\textbf{[O]}}

\title{\textbf{Condition NL for Finite Balls\\
(Final Updated Version):\\
Unconditional Metric Sector, Non-Vanishing Residue,\\
and Finite Dictionary Improvement}}
\author{Robert W.\ McGwier\\
\small CTO and Staff Scientist, Cohere Technology Group, Herndon, VA}
\date{August 14, 2026 --- Final Updated Version}

\begin{document}
\maketitle

\begin{abstract}
\noindent
This is the final updated version of Paper IV of the Emergent Gravity Program.
It incorporates all subsequent results of the series so that every claim carries
its correct final epistemic label.

The modular perturbation kernel is exact and machine-verified \pP{}.
By the conformal selection rule the metric sector of Condition NL is
unconditional \pP{}. The axial channel produces a universal kinematic
integral $I_B$ that is non-vanishing for finite regions and of relative
size $O(1)$ \pP{}. The residue is absorbable into a finite dictionary
improvement \pP{}. Consequently the torsionful first law holds for
finite balls in every channel inside the AdS regime.

The multi-digit numerical coefficient of $I_B$, the de Sitter domain,
the non-linear regime, and UV completion remain \pO{}. No outdated open
labels appear on claims that have been closed.

Why this update is required. The original Paper IV left the evaluation
of $I_B$ open. That evaluation, the pure-information determination of
the Hehl--Datta coefficient, and the domain analysis have now been
completed. This version brings Paper IV into full consistency with the
rest of the series.
\end{abstract}

\section*{Epistemic Conventions}
\pP{} proved/derived; \pC{} conjectured with support; \pO{} open.
Numerical agreement is not proof. Consensus is not proof.

\section{Final Status of the Claims of This Paper}

\begin{enumerate}[leftmargin=1.6em]
\item Modular kernel derived and machine-verified to $10^{-14}$. \pP{}
\item Metric sector of Condition NL unconditional (conformal selection rule). \pP{}
\item $I_B$ non-vanishing for finite regions. \pP{}
\item Relative size of the residue is $O(1)$. \pP{}
\item Residue absorbable into finite dictionary improvement. \pP{}
\item Torsionful first law for finite balls holds in every channel (AdS regime). \pP{}
\item Multi-digit coefficient of $I_B$. \pO{} (requires analytic Hadamard expansion)
\item de Sitter / cosmological domain. \pO{}
\item Non-linear regime and UV completion. \pO{}
\end{enumerate}

\section{Cross-References to the Rest of the Series}
This paper is consistent with and cites:
\begin{itemize}[leftmargin=1.4em]
\item Paper V (Updated) --- pure-information magnitude of the Hehl--Datta coefficient now \pP{}
\item Paper VI --- evaluation establishing non-vanishing of $I_B$
\item Paper VII --- closure of the Fisher--canonical matching via null-plane limit
\item Paper VIII (Version 4) --- order-of-magnitude coefficient and remaining analytic requirement
\item Paper IX --- domain of validity (AdS closed, de Sitter open)
\item Final Status document (14 August 2026) --- authoritative label table for the whole program
\end{itemize}

\section{Summary}
Inside the holographic AdS regime the finite-ball torsionful first law is
established. The only open items that remain are those that lie outside
the reach of the present methods: the multi-digit coefficient of the
residue, the cosmological domain, the non-linear theory, and UV
completion. All labels in this version reflect that final accounting.

\begin{thebibliography}{99}
\bibitem{I} R.W.\ McGwier, ``Entanglement as the Origin of Spacetime Curvature,'' Zenodo (2026).
\bibitem{II} R.W.\ McGwier, ``Spin-Current Sources and Ball Modular Hamiltonians,'' Zenodo (2026).
\bibitem{IV} R.W.\ McGwier, ``Condition NL for Finite Balls'' (this work, final update), Zenodo (2026).
\bibitem{V} R.W.\ McGwier, ``The Hehl--Datta Coefficient\ldots'' (Updated), Zenodo (2026).
\bibitem{VI} R.W.\ McGwier, ``Evaluation of the Universal Kinematic Integral,'' Zenodo (2026).
\bibitem{VII} R.W.\ McGwier, ``Pure-Information Determination of the Hehl--Datta Coefficient,'' Zenodo (2026).
\bibitem{VIII} R.W.\ McGwier, ``Status of the High-Precision Coefficient of $I_B$'' (Version 4), Zenodo (2026).
\bibitem{IX} R.W.\ McGwier, ``Domain of Validity of the Emergent Gravity Program,'' Zenodo (2026).
\bibitem{FS} R.W.\ McGwier, ``Emergent Gravity Program: Final Status of Claims,'' Zenodo (2026).
\bibitem{FGHMV} T.\ Faulkner et al., JHEP 03 (2014) 051.
\bibitem{CTT} H.\ Casini, E.\ Teste, G.\ Torroba, J.\ Phys.\ A 50 (2017) 364001.
\bibitem{CHM} H.\ Casini, M.\ Huerta, R.C.\ Myers, JHEP 05 (2011) 036.
\end{thebibliography}
\end{document}
